% ============================================================================
% preamble.tex -- shared packages, macros and environments
%
% Deliberately template-agnostic: no \documentclass, no page geometry beyond a
% sane default, no font selection. When the university template arrives, delete
% the "Layout" block below and keep everything else.
% ============================================================================

% --- Layout (replace with the institutional template) -----------------------
\usepackage[a4paper,margin=3cm]{geometry}
\usepackage{microtype}
\usepackage{parskip}

% --- Mathematics ------------------------------------------------------------
\usepackage{amsmath}
\usepackage{amssymb}
\usepackage{amsthm}
\usepackage{mathtools}
\usepackage{bm}
\usepackage{siunitx}
\sisetup{
  detect-all,
  range-units        = single,
  separate-uncertainty = true,
  group-minimum-digits = 4,
}

% --- Tables, floats, graphics ----------------------------------------------
\usepackage{booktabs}
\usepackage{multirow}
\usepackage{array}
\usepackage{graphicx}
\usepackage{float}
\usepackage{longtable}

% --- Code -------------------------------------------------------------------
\usepackage{listings}
\usepackage{xcolor}

\definecolor{codebg}{HTML}{F7F7F9}
\definecolor{coderule}{HTML}{D8D8DE}
\definecolor{codekw}{HTML}{0B5FA5}
\definecolor{codestr}{HTML}{9A2617}
\definecolor{codecom}{HTML}{6A737D}
\definecolor{accent}{HTML}{1B4F72}
\definecolor{warnbg}{HTML}{FDF6E3}
\definecolor{warnrule}{HTML}{B58900}

\lstdefinestyle{pystyle}{
  language          = Python,
  basicstyle        = \ttfamily\footnotesize,
  keywordstyle      = \color{codekw}\bfseries,
  stringstyle       = \color{codestr},
  commentstyle      = \color{codecom}\itshape,
  numberstyle       = \tiny\color{codecom},
  backgroundcolor   = \color{codebg},
  frame             = single,
  rulecolor         = \color{coderule},
  framesep          = 5pt,
  xleftmargin       = 10pt,
  xrightmargin      = 4pt,
  breaklines        = true,
  breakatwhitespace = true,
  showstringspaces  = false,
  columns           = fullflexible,
  keepspaces        = true,
  literate          = {-}{{-}}1 {*}{{*}}1,
}
\lstset{style=pystyle}

% --- Callout boxes ----------------------------------------------------------
\usepackage[most]{tcolorbox}
\tcbuselibrary{breakable}

% "What the code actually does" -- used wherever the implementation departs
% from the idealised statement in the surrounding text.
\newtcolorbox{implnote}[1][]{
  breakable,
  enhanced,
  colback   = codebg,
  colframe  = accent,
  boxrule   = 0.4pt,
  leftrule  = 2.5pt,
  arc       = 1pt,
  fonttitle = \bfseries\sffamily\small,
  title     = {Implementation: #1},
  top = 4pt, bottom = 4pt,
}

% Measured-fact box: numbers that came out of a run, not out of an argument.
\newtcolorbox{measured}[1][]{
  breakable,
  enhanced,
  colback   = white,
  colframe  = accent!70,
  boxrule   = 0.4pt,
  arc       = 1pt,
  fonttitle = \bfseries\sffamily\small,
  title     = {Measured: #1},
  top = 4pt, bottom = 4pt,
}

% Caveat box: a limitation of the method or of the present validation state.
\newtcolorbox{caveat}[1][]{
  breakable,
  enhanced,
  colback   = warnbg,
  colframe  = warnrule,
  boxrule   = 0.4pt,
  leftrule  = 2.5pt,
  arc       = 1pt,
  fonttitle = \bfseries\sffamily\small,
  title     = {Caveat: #1},
  top = 4pt, bottom = 4pt,
}

% --- Theorem-like environments ---------------------------------------------
\theoremstyle{definition}
\newtheorem{definition}{Definition}[chapter]
\newtheorem{convention}[definition]{Convention}
\newtheorem{example}[definition]{Example}

\theoremstyle{plain}
\newtheorem{proposition}[definition]{Proposition}
\newtheorem{lemma}[definition]{Lemma}

\theoremstyle{remark}
\newtheorem{remark}[definition]{Remark}

% --- Cross-referencing ------------------------------------------------------
\usepackage[colorlinks=true,linkcolor=accent,citecolor=accent,urlcolor=accent]{hyperref}
\usepackage[capitalise,noabbrev]{cleveref}

% ============================================================================
% Notation macros
% ============================================================================

% Reduced flux quantum \varphi_0 = \Phi_0 / 2\pi
\newcommand{\rphi}{\varphi_{0}}
\newcommand{\fluxquantum}{\Phi_{0}}

% Operators and sets
\DeclareMathOperator{\Real}{Re}
\DeclareMathOperator{\Imag}{Im}
\DeclareMathOperator{\diag}{diag}
\DeclareMathOperator{\blockdiag}{blockdiag}
\DeclareMathOperator*{\argmin}{arg\,min}
\DeclareMathOperator{\spanof}{span}
\DeclareMathOperator{\sgn}{sgn}
\DeclareMathOperator{\nnz}{nnz}

% Matrices (upright bold)
\newcommand{\mat}[1]{\mathbf{#1}}
\newcommand{\vect}[1]{\mathbf{#1}}
\newcommand{\Cmat}{\mat{C}}
\newcommand{\Gmat}{\mat{G}}
\newcommand{\Kmat}{\mat{K}}
\newcommand{\Bphi}{\mat{B}_{\phi}}
\newcommand{\Dmat}{\mat{D}}
\newcommand{\Jmat}{\mat{J}}
\newcommand{\Mmat}{\mat{M}}
\newcommand{\Smat}{\mat{S}}
\newcommand{\Imat}{\mat{I}}
\newcommand{\Khat}{\widehat{\mat{K}}}
\newcommand{\ghat}{\hat{\gamma}}

% Vectors of unknowns
\newcommand{\Xvec}{\vect{X}}
\newcommand{\Rvec}{\vect{R}}
\newcommand{\Nvec}{\vect{N}}
\newcommand{\Svec}{\vect{S}}
\newcommand{\Vvec}{\vect{V}}
\newcommand{\xvec}{\vect{x}}
\newcommand{\isrc}{\vect{i}_{\mathrm{src}}}
\newcommand{\iJ}{\vect{i}_{J}}
\newcommand{\Ic}{\vect{I}_{c}}

% Frequencies.
% \wp is LaTeX's Weierstrass-p symbol, which this document never uses; it is
% deliberately repurposed as the pump angular frequency.
\renewcommand{\wp}{\omega_{p}}
\newcommand{\ws}{\omega_{s}}
\newcommand{\wi}{\omega_{i}}

% Mode sets
\newcommand{\modeset}{\mathcal{K}}
\newcommand{\toneset}{\mathcal{T}}
\newcommand{\krylov}{\mathcal{K}}

% Miscellaneous
\newcommand{\code}[1]{\texttt{\small #1}}
\newcommand{\file}[1]{\texttt{\small #1}}

% Logarithmic units. Declared only as siunitx units -- do not also define
% \dB/\dBm as plain macros, or the declaration and the macro collide.
\DeclareSIUnit{\dBm}{dBm}
\DeclareSIUnit{\dB}{dB}

% Elementwise (Hadamard) product
\newcommand{\hada}{\odot}
